\documentclass[10pt,a4paper]{amsart}
\usepackage[margin=19mm]{geometry}
\usepackage{amsmath,amssymb,mathtools}
\usepackage[scr=rsfs,cal=euler]{mathalfa}
\usepackage{xfrac}
\usepackage[unicode,hidelinks,bookmarks=false]{hyperref}
\newcommand{\ie}{i.e.\ }
\newcommand{\cf}{cf.\ }
\newcommand{\resp}{respectively\ }
\newcommand{\Subsection}{\subsection}
\providecommand{\nobreakdash}{\nobreak\hbox{-}}
\setlength{\emergencystretch}{2em}
\multlinegap0pt
\usepackage{bm}
\usepackage{bbm}
\usepackage{dsfont}
\usepackage{xspace}
\usepackage{cancel}
\usepackage{mathtools}
\usepackage{amsthm}
\newtheorem{lemma}{Lemma}
\renewcommand{\pmod}[1]{\mkern3mu(\mathrm{mod}\ #1)}
\title[On the hyperbolic prime number theorem]{On the hyperbolic prime number theorem}
\author{Alisa Sedunova}
\date{Source: arXiv:2609.28200v1 (CC BY 4.0)}
\begin{document}
\maketitle

\noindent\emph{Source and licence.} On the hyperbolic prime number theorem. Version arXiv:2609.28200v1 (CC BY 4.0), distributed under the Creative Commons Attribution 4.0 International licence, as recorded in the arXiv licence metadata for this exact version and shown on the abstract page https://arxiv.org/abs/2609.28200v1. Full rights notice: licenses/arxiv-2609.28200-CC-BY-4.0.txt.\par\smallskip

\noindent\textit{Editorial scope.} Counted unit: the Lemma whose source label is \texttt{convr-even} in the article (source0.tex lines 373--381, source label \texttt{convr-even}), delivered together with the article's own complete original proof of it (source0.tex lines 383--488, one \texttt{proof} environment of 106 lines). No mathematics is written, repaired or re-derived: statement and proof are the article's own text line for line. Required context retained uncounted: rinAps (lines 198--205); convr (lines 207--218). Every definition, display and label of those lines is retained unchanged. Omitted from this package and not redistributed: 1--197, 206--206, 219--372, 382--382, 489--930. Nothing inside a retained excerpt is omitted or altered: every retained body is delivered exactly as the article's lines are written, so no omission receipt of any kind accompanies this package. Citations: the retained excerpts carry no citation command at all, so no bibliography entry of the article is transcribed and no reference is redistributed. Reference adaptation: none was needed and none was made; every reference inside the delivered unit resolves inside it, so no unresolved reference or citation is produced. Printed slips kept literal: no printed word, symbol, hypothesis or number of the counted statement or of its proof was altered. Layout and class: the article's own class is \texttt{amsart} with packages including dsfont, amsmath, amsthm, amssymb, multirow, mathtools, indentfirst, graphicx, color, hyperref, comment; the wrapper is the lane's pinned \texttt{amsart} editorial wrapper at 10pt on A4, which restates the theorem-like environments the retained text needs and adds the article's own macro definitions that the retained text needs (\texttt{\textbackslash pmod}), with extra wrapper packages (bm, bbm, dsfont, xspace, cancel, mathtools). The delivered numbering is the unit's own. Assets: no retained span contains a figure environment, a graphics-inclusion command, a PDF-page inclusion, a table or a TikZ picture; no image file is copied into this package, the package declares no asset dependency and no image file is redistributed with it. Licence: version arXiv:2609.28200v1 of this article is distributed under the Creative Commons Attribution 4.0 International licence, as recorded in the arXiv licence metadata for this exact version and shown on the abstract page https://arxiv.org/abs/2609.28200v1. A full rights notice is delivered at licenses/arxiv-2609.28200-CC-BY-4.0.txt. Characters: every retained excerpt is pure ASCII, so the delivered engine, XeLaTeX with its default Latin Modern text font, reports no missing character.
\begin{lemma}\label{rinAps}
	Let $(a,q)=1$. 
	Then for $q \not\equiv 0 \pmod 4$ we have
	\[
		\sum_{\substack{n \leq x \\ n \equiv a \pmod q}} r(n) = \frac{\pi x}{q}\prod_{p \mid q} \left(1-\frac{\chi_4(p)}{p}\right) + O (q^{-1/2}x^{2/3+\varepsilon})
	\]		
and the main term is multiplied by $1+\chi_4(a)$ for $q \equiv 0 \pmod 4$.
\end{lemma} 

\begin{lemma} \label{convr}
    Let $d$ be an odd positive integer and let $\varepsilon > 0$. Then
    \[
        \sum_{\substack{n\le x\\ n\equiv3\ (\mathrm{mod}\ 4)\\ n\equiv0\ (\mathrm{mod}\ d)}}r(n-2)r(n+2) = 8g(d)x+O\big(d^{-1/2}x^{11/12+\varepsilon}+x^{1/2+\varepsilon}\big),
    \]
    where
    \[
        g(d) = \frac{1}{d} \prod_{p\mid d}\frac{p-\chi_4(p)}{p+\chi_4(p)}
    \]
    and the implied constant depends only on $\varepsilon$. 
    In particular the error term is $O(x^{11/12+\varepsilon})$ uniformly in $d$.
\end{lemma}

\begin{lemma}\label{convr-even}
    Let $q$ be an odd positive integer, let $a$ be an integer with $(a(a+1),q) = 1$ and let
    $\varepsilon > 0$. Then
    \[
        \sum_{\substack{n \le x\\ n \equiv a \pmod q}} r(n) r(n+1)
        = 8g(q)x + O\bigl(q^{-1/2}x^{11/12+\varepsilon}+x^{1/2+\varepsilon}\bigr),
    \]
    with $g$ as in Lemma \ref{convr} and the implied constant depending only on $\varepsilon$.
\end{lemma}

\begin{proof}
    Exactly one of $n$, $n+1$ is odd, and for odd $m$ one has $r(m) = 0$ unless $m \equiv 1 \pmod 4$. If $n \equiv 2 \pmod 4$ then $n+1 \equiv 3 \pmod 4$, so $r(n+1) = 0$.
    Similarly, if $n \equiv 3 \pmod 4$ then $r(n) = 0$. 
    Hence only $n \equiv 0, 1 \pmod 4$ contribute to the sum.
    We work these cases out separately.

    First, let $n \equiv 0 \pmod 4$ and denote its contribution by $S_1(x)$. As $n + 1\equiv 1 \pmod 4$ we have
    \[
        r(n+1) = 8\sum_{\substack{k \mid n+1\\ k < \sqrt{n + 1}}}\chi_4(k) + 4\chi_4(\sqrt{n+1}),
    \]
    and
    \[
        \begin{split}
            S_1(x) &= 8 \sum_{\substack{n \le x\\ n \equiv a \pmod q \\ n \equiv 0 \pmod 4}} r(n) \sum_{\substack{k \mid n+1\\ k < \sqrt{n + 1}}}\chi_4(k) + O(x^{1/2 + \varepsilon})\\
            &= 8 \sum_{k < \sqrt{x+1}} \chi_4(k) \sum_{\substack{n \le x \\ n \equiv 0 \pmod 4\\ n \equiv a \pmod q \\ n \equiv -1 \pmod k\\ n > k^2-1}} r(n) + O(x^{1/2 + \varepsilon}).
        \end{split}
    \]
    Here $k$ is odd and $(k,q) = 1$ --- otherwise the inner sum is empty, since a prime $p \mid (k,q)$ would divide $n+1$ and $q$, hence $a+1$, contrary to hypothesis. Writing $k \mid n+1$ as $n \equiv -1 \pmod k$ and substituting $n = 4m$, so that $r(n) = r(m)$, the conditions on $m$ become a single class $a' \pmod{qk}$ coprime to $qk$, and $(k^2-1)/4 < m \le x/4$ (note that $(k^2-1)/4$ is an integer as $k$ is odd). Since $qk$ is odd, Lemma \ref{rinAps} applies
    \[
        \begin{split}
            \sum_{\substack{n \le x \\ n \equiv 0 \pmod 4\\ n \equiv a \pmod q \\
                n \equiv -1 \pmod k\\ n > k^2-1}} r(n)
            &= \sum_{\substack{m \le x/4 \\ m \equiv a' \pmod {qk} \\
                m > (k^2-1)/4}} r(m)\\
            &= \frac{\pi (x-k^2)}{4qk}h(q)h(k)
               + O\bigl((qk)^{-1/2}(x^{2/3+\varepsilon}+k^{4/3+\varepsilon})\bigr),
        \end{split}
    \]
    the $+1$ left over from $x-k^2+1$ being absorbed into the error term.
    
    Summing the error over $k$ we obtain
    \[
        \begin{split}
            \sum_{k < \sqrt{x+1}}(qk)^{-1/2}\bigl(x^{2/3+\varepsilon}+k^{4/3+\varepsilon}\bigr)
            &\ll q^{-1/2}x^{2/3+\varepsilon}\sum_{k \ll \sqrt{x}} k^{-1/2} \\
            &\ll q^{-1/2} x^{11/12+\varepsilon}.
        \end{split}
    \]
    In the main term $\chi_4$ vanishes at even $k$, thus
    \[
        8\sum_{\substack{k<\sqrt{x+1}\\ (k,q)=1}}\chi_4(k)\frac{\pi(x-k^2)}{4qk}h(q)h(k)
        = \frac{2\pi h(q)}{q}\sum_{\substack{k<\sqrt{x+1}\\ (k,q)=1}}
          \chi_4(k)h(k)\Bigl(\frac{x}{k}-k\Bigr).
    \]
    As in the proof of Lemma \ref{convr} we have
    \[
        \sum_{\substack{k \le t \\ (k,q)=1}} \chi_4(k)h(k) \ll \tau(q)\log t,
    \]
    whence by partial summation
    \[
        \sum_{\substack{k<\sqrt{x+1}\\ (k,q)=1}}\chi_4(k)h(k)k \ll \tau(q)x^{1/2}\log x,
        \quad
        \sum_{\substack{k \ge \sqrt{x+1}\\ (k,q)=1}}\frac{\chi_4(k)h(k)}{k}
        \ll \frac{\tau(q)\log x}{x^{1/2}}.
    \]
    Since $h(q)\tau(q) \ll q^{\varepsilon}$, both contribute $O(x^{1/2+\varepsilon})$
    after multiplication by $2\pi h(q)/q$ and $2\pi h(q)x/q$ respectively. On completing the sum to infinity we get
    \[
        \sum_{\substack{k \ge 1\\ (k,q)=1}}\frac{\chi_4(k)h(k)}{k}
        = \prod_{p \nmid q}\Bigl(1+\frac{\chi_4(p)}{p}\Bigr)
        = \frac{2}{\pi}\prod_{p \mid q}\Bigl(1+\frac{\chi_4(p)}{p}\Bigr)^{-1},
    \]
    so that the main term equals
    \[
        \frac{2\pi h(q)x}{q}\cdot\frac{2}{\pi}
        \prod_{p \mid q}\Bigl(1+\frac{\chi_4(p)}{p}\Bigr)^{-1}
        = \frac{4x}{q}\prod_{p \mid q}\frac{p-\chi_4(p)}{p+\chi_4(p)} = 4g(q)x.
    \]
    Therefore
    \[
        S_1(x) = 4g(q)x + O\bigl(q^{-1/2}x^{11/12+\varepsilon}+x^{1/2+\varepsilon}\bigr).
    \]
    
    Next, let $n \equiv 1 \pmod 4$ and denote its contribution by $S_2(x)$. This
    time $n$ itself is odd and $n \equiv 1 \pmod 4$, so
    \[
        r(n) = 8\sum_{\substack{k \mid n\\ k < \sqrt{n}}}\chi_4(k) + 4\chi_4(\sqrt{n}),
    \]
    therefore
    \[
        S_2(x) = 8\sum_{k<\sqrt{x}}\chi_4(k)
        \sum_{\substack{n \le x \\ n \equiv 1 \pmod 4\\ n \equiv a \pmod q \\
            n \equiv 0 \pmod k\\ n > k^2}} r(n+1) + O(x^{1/2+\varepsilon}).
    \]
    Again $k$ is odd and $(k,q)=1$, since a prime $p \mid (k,q)$ would divide $n$
    and $q$, hence $a$. Substituting $n+1 = 2m$, so that $r(n+1)=r(m)$, the
    condition $n \equiv 1 \pmod 4$ makes $m$ odd, the conditions on $m$ become a
    single class $a'' \pmod{2qk}$ coprime to $2qk$, and
    $(k^2+1)/2 < m \le (x+1)/2$. As $\chi_4(2)=0$ we have $h(2qk)=h(q)h(k)$, and
    since $4 \nmid 2qk$ Lemma \ref{rinAps} gives
    \[
        \begin{split}
            \sum_{\substack{m \le (x+1)/2 \\ m \equiv a'' \pmod{2qk} \\
                m > (k^2+1)/2}}& r(m)
            = \frac{\pi}{2qk}\Bigl(\frac{x+1}{2}-\frac{k^2+1}{2}\Bigr)h(q)h(k) \\
               &+ O\bigl((qk)^{-1/2}(x^{2/3+\varepsilon}+k^{4/3+\varepsilon})\bigr)\\
            &= \frac{\pi (x-k^2)}{4qk}h(q)h(k)
               + O\bigl((qk)^{-1/2}(x^{2/3+\varepsilon}+k^{4/3+\varepsilon})\bigr).
        \end{split}
    \]
    Proceeding as in the case of $S_1(x)$ we conclude
    \[
        S_2(x) = 4g(q)x + O\bigl(q^{-1/2}x^{11/12+\varepsilon}+x^{1/2+\varepsilon}\bigr).
    \]
    Adding $S_1(x)$ and $S_2(x)$ completes the proof.
\end{proof}

\end{document}
